% Source: https://artofproblemsolving.com/community/c3h2013862_question_on_a_tetrahedron_problem_someone_posted_a_few_days_ago
% Stable identity: aops-topic-2013862-q01
\textbf{(a)} The tetrahedron $T$ is a pyramid with apex $C$ and base the right triangle $OAB$ in the plane $z=0$; this base consists of the points $(u,v,0)$ with $u\ge0$, $v\ge0$, $u+v\le1$. The points of $T$ have heights $0\le z\le1$, and the cross-section of the pyramid at height $h$ is the image of the base under the dilation with centre $C$ and ratio $1-h$. This dilation sends $(u,v,0)$ to
\[C+(1-h)\bigl((u,v,0)-C\bigr)=\bigl((1-h)u,\ (1-h)v,\ h\bigr).\]
So the cross-section at height $h$ consists of the points $(x,y,h)$ with $x=(1-h)u$, $y=(1-h)v$, that is, with
\[x\ge0,\qquad y\ge0,\qquad x+y\le1-h .\]

\begin{asy}
usepackage("amsmath");
size(5.5cm);
defaultpen(fontsize(10pt));
pair pr(real x,real y,real z){return (x+0.45*y,z+0.35*y);}
pair O=pr(0,0,0), A=pr(1,0,0), B=pr(0,1,0), C=pr(0,0,1);
real h=0.4;
pair P1=pr(0,0,h), P2=pr(1-h,0,h), P3=pr(0,1-h,h);
fill(A--B--C--cycle,paleblue);
draw(pr(0,1,0)--pr(1,1,0),gray+dashed); draw(pr(0,1,0)--pr(0,1,1),gray+dashed);
draw(A--pr(1,0,1)--C,gray); draw(A--pr(1,1,0)--pr(1,1,1)--pr(1,0,1),gray); draw(C--pr(0,1,1)--pr(1,1,1),gray);
fill(P1--P2--P3--cycle,palered);
draw(P1--P2,red+linewidth(0.8)); draw(P2--P3--P1,red+dashed);
draw(O--A,blue+linewidth(0.9)); draw(O--C,blue+linewidth(0.9)); draw(A--C,blue+linewidth(0.9));
draw(O--B,blue+dashed); draw(A--B,blue+dashed); draw(B--C,blue+dashed);
dot(O); dot(A); dot(B); dot(C);
label("$h$",pr(0,0,h/2),W,red);
label("$O$",O,SW); label("$A(1,0,0)$",A,S); label("$B$",B,S); label("$C(0,0,1)$",C,W);
label("\mbox{}",(min(currentpicture,user=true)+max(currentpicture,user=true))/2,black); // keeps the text after the figure black when LaTeX typesets the labels inline
\end{asy}

Hence $(x,y,z)$ lies in $T$ exactly when $0\le z\le1$, $x\ge0$, $y\ge0$ and $x+y\le1-z$, which is the same as $x\ge0$, $y\ge0$, $z\ge0$ and $x+y+z\le1$ (the condition $z\le1$ then follows automatically). $\blacksquare$

The face $ABC$ is the part of $T$ in the plane $x+y+z=1$ (indeed $A$, $B$, $C$ satisfy $x+y+z=1$), and the other three faces lie in the planes $x=0$, $y=0$, $z=0$. A point inside $T$ is on none of the faces, so all four inequalities are strict there; in particular $x+y+z<1$.

\textbf{(b)} Every point with $x,y,z\ge0$ and $x+y+z\le1$ has each coordinate at most $1$, so $T$ lies inside the unit cube. By (a), the points of the cube with $x+y+z<1$ are the points of $T$ not on the face $ABC$, and that face has volume $0$. Since the cube has volume $1$, the probability is the volume of $T$, a pyramid with base area $[OAB]=\dfrac{1}{2}$ and height $OC=1$:
\[\dfrac{1}{3}\cdot\dfrac{1}{2}\cdot1=\boxed{\dfrac{1}{6}}.\]
